% =====================================================================
%  3.5 Eksponential- og potensfunksjoner — teori og oppgaver
% =====================================================================

\del{1}{Eksponentialfunksjoner}

\begin{ordliste}
eksponential\-funksjon & En funksjon $f(x) = a \cdot b^x$ med $a > 0$ og $b > 0$. & \mittsprak\\
startverdi & Tallet $a$. Det er verdien når $x = 0$, fordi $f(0) = a$. & \mittsprak\\
vekstfaktor & Tallet $b$. Verdien ganges med $b$ for hver gang $x$ øker med 1. & \mittsprak\\
doblingstid & Tiden det tar før verdien er doblet. & \mittsprak\\
halveringstid & Tiden det tar før verdien er halvert. & \mittsprak\\
\end{ordliste}

\begin{mustkunne}
\begin{minipage}[t]{0.50\linewidth}
\textbf{Lineær: like store tillegg. Eksponentiell: samme faktor.}\par\smallskip
\begin{tabular}[t]{@{}c|cccc@{}}
$x$ & 0 & 1 & 2 & 3\\ \hline
$100 + 10x$ & 100 & 110 & 120 & 130\\
$100 \cdot 1{,}1^x$ & 100 & 110 & 121 & 133,1\\
\end{tabular}\par\smallskip
Øverst legger vi til 10 hver gang. Nederst ganger vi med $1{,}1$ hver gang. Derfor er $f(x) = a \cdot b^x$: start på $a$, og gang med $b$ for hvert steg.
\end{minipage}\hfill
\begin{minipage}[t]{0.47\linewidth}
\vspace{0pt}\centering
\begin{tikzpicture}[baseline=(current bounding box.north)]
\begin{axis}[width=\linewidth, height=0.72\linewidth, axis lines=middle, xmin=-3.5, xmax=3.5, ymin=-0.5, ymax=8.8,
  xtick={-2,0,2}, ytick={2,4,6,8}, tick label style={font=\scriptsize}, xlabel={$x$}, ylabel={$y$}, label style={font=\small}]
\addplot[domain=-3.3:3, thick, hoved, samples=50] {2^x};
\addplot[domain=-3:3.3, thick, red!70!black, samples=50] {(1/2)^x};
\node[font=\scriptsize, hoved, anchor=west] at (axis cs:1.2,7) {$2^x$};
\node[font=\scriptsize, red!70!black, anchor=east] at (axis cs:-1.2,7) {$\left(\tfrac12\right)^x$};
\end{axis}
\end{tikzpicture}\par
{\footnotesize $b > 1$: voksende. $0 < b < 1$: avtagende.}
\end{minipage}\medskip

\textbf{Vekstfaktor fra prosent:} vekst på $p\,\%$ gir $b = 1 + \frac{p}{100}$. Nedgang på $p\,\%$ gir $b = 1 - \frac{p}{100}$.
Grafen har $y = 0$ som horisontal asymptote og skjærer $y$-aksen i $(0, a)$. Verdien blir aldri negativ.
\end{mustkunne}

\begin{eksempel}[1 — rente på rente]
Tenkt eksempel: 20\,000 kr står på en konto med 4\,\% rente per år. Finn beløpet etter 5 år.

\begin{stegtabell}
$b = 1 + \frac{4}{100} = 1{,}04$ & Vekst på 4\,\% gir vekstfaktor $1{,}04$.\\
$f(x) = 20\,000 \cdot 1{,}04^x$ & $a = 20\,000$ er startbeløpet. $x$ er antall år.\\
$f(5) = 20\,000 \cdot 1{,}04^5 \approx 24\,333$ & Kalkulator. Beløpet er omtrent 24\,333 kr.\\
Tolking: & $a$ er beløpet i dag. $b$ sier at beløpet ganges med $1{,}04$ hvert år.\\
\end{stegtabell}
\end{eksempel}

\begin{eksempel}[2 — finn funksjonen fra to punkter]
Finn $f(x) = a \cdot b^x$ som går gjennom $(2, 12)$ og $(5, 96)$.

\begin{stegtabell}
$a \cdot b^2 = 12$ og $a \cdot b^5 = 96$ & Sett inn begge punktene.\\
$\dfrac{a \cdot b^5}{a \cdot b^2} = \dfrac{96}{12}$, altså $b^3 = 8$ & Del den andre likningen på den første. Da forsvinner $a$.\\
$b = 2$ & Fordi $2^3 = 8$.\\
$a \cdot 2^2 = 12$, så $a = 3$ & Sett $b$ inn i den første likningen.\\
Kontroll: $3 \cdot 2^2 = 12$ og $3 \cdot 2^5 = 96$ & Begge punktene stemmer.\\
\end{stegtabell}
\end{eksempel}

\begin{vanligfeil}
\begin{itemize}[leftmargin=1.2em,itemsep=1pt,topsep=0pt]
  \item «5\,\% vekst gir $b = 0{,}05$ eller $b = 1{,}5$» er feil. Riktig er $b = 1 + \frac{5}{100} = 1{,}05$.
  \item «20\,\% nedgang gir $b = 0{,}2$» er feil. Riktig er $b = 1 - \frac{20}{100} = 0{,}8$.
  \item «5\,\% per år i 10 år er 50\,\%» er feil. Det blir $1{,}05^{10} \approx 1{,}63$, altså omtrent 63\,\% vekst.
  \item $a \cdot b^x = (ab)^x$ gjelder ikke generelt. For $2 \cdot 5^x$ og $x = 2$ er $2 \cdot 5^2 = 50$, men $(2 \cdot 5)^2 = 100$.
\end{itemize}
\end{vanligfeil}

\nivaahode{1}

\oppg{1.1}{1}{\uh}
Finn vekstfaktoren.
\begin{tasks}(2)
  \task Verdien øker med $12\,\%$.
  \task Verdien synker med $35\,\%$.
  \task Verdien øker med $2{,}5\,\%$.
  \task Verdien øker med $60\,\%$.
\end{tasks}

\oppg{1.2}{1}{\uh}
Se på $f(x) = 4 \cdot 3^x$.
\begin{tasks}(1)
  \task Hva er startverdien og vekstfaktoren?
  \task Regn ut $f(0)$, $f(1)$ og $f(2)$.
  \task Er funksjonen voksende eller avtagende?
\end{tasks}

\oppg{1.3}{1}{\uh}
Tenkt eksempel: En vare koster 500 kr. Prisen øker med $10\,\%$ per år. Fullfør.
\begin{tasks}(1)
  \task Vekstfaktoren er $b = 1 + \dfrac{10}{100} = $ \skrivlinje[2cm]
  \task Funksjonen er $f(x) = 500 \cdot$ \skrivlinje[2cm]$^x$
  \task Prisen etter 2 år er $f(2) = $ \skrivlinje[2cm]
\end{tasks}

\nivaahode{2}

\oppg{1.4}{2}{\mh}
Tenkt eksempel: En bil er verdt 240\,000 kr. Verdien synker med $15\,\%$ per år.
\begin{tasks}(1)
  \task Skriv en funksjon $f$ som gir verdien etter $x$ år.
  \task Finn verdien etter 4 år. Rund av til nærmeste hundre kroner.
  \task Når er verdien 100\,000 kr? Bruk graf eller digitalt verktøy. Svar med én desimal.
  \task Modellen gir at bilen er verdt ca. 1\,800 kr etter 30 år. Drøft om modellen er gyldig.
\end{tasks}

\oppg{1.5}{2}{\uh}
Finn $f(x) = a \cdot b^x$ som går gjennom $(1, 40)$ og $(3, 10)$. Kontroller svaret.

\oppg{1.6}{2}{\uh}
Tenkt eksempel: En vare koster 800 kr og blir $25\,\%$ billigere hvert år. Hva koster den etter 2 år?
Finn \textbf{den første feilen}. Regn riktig.
\[ b = 1 - 0{,}25 = 0{,}75 \ \longrightarrow\ f(x) = 800 \cdot 0{,}75^x \ \longrightarrow\ f(2) = (800 \cdot 0{,}75)^2 \ \longrightarrow\ f(2) = 600^2 = 360\,000 \]

\oppg{1.7}{2}{\mh}
Tenkt eksempel: Et beløp vokser med $6\,\%$ per år.
\begin{tasks}(1)
  \task Startbeløpet er 1\,000 kr. Finn doblingstiden grafisk eller digitalt. Svar med én desimal.
  \task Finn doblingstiden når startbeløpet er 3\,000 kr. Hva legger du merke til?
\end{tasks}

\oppg{1.8}{2}{\mh}
Tenkt eksempel: To månedslønner utvikler seg slik. $x$ er antall år.
\[ A(x) = 25\,000 + 2\,000x \qquad B(x) = 25\,000 \cdot 1{,}08^x \]
\begin{tasks}(1)
  \task Regn ut $A$ og $B$ for $x = 1$, $x = 2$ og $x = 10$.
  \task Når er $B$ større enn $A$ for første gang?
  \task Forklar med egne ord forskjellen på hvordan de to lønnene vokser.
\end{tasks}

\nivaahode{3}

\oppg{1.9}{3}{\uh}
Er påstanden alltid, noen ganger eller aldri sann? Begrunn svaret.
Vi ser på $f(x) = a \cdot b^x$ og $g(x) = c \cdot d^x$ med positive $a$, $b$, $c$ og $d$.
\begin{tasks}(1)
  \task Produktet $f(x) \cdot g(x)$ er en eksponentialfunksjon.
  \task Summen $f(x) + g(x)$ er en eksponentialfunksjon.
  \task Vekst på $30\,\%$ to år på rad gir $60\,\%$ vekst til sammen.
\end{tasks}

\oppg{1.10}{3}{\uh}
Verdien til en eksponentialfunksjon $f(x) = a \cdot b^x$ blir doblet hvert 3. år.
\begin{tasks}(1)
  \task Forklar hvorfor $b^3 = 2$.
  \task Hvor mange ganger større er verdien etter 9 år? Forklar uten å finne $b$.
\end{tasks}

\oppg{1.11}{3}{\uh}
\begin{tasks}(1)
  \task Lag selv en avtagende eksponentialfunksjon $f$ med $f(3) = 6$.
  \task En eksponentialfunksjon $f(x) = a \cdot b^x$ har $f(0) = 5$ og $f(2) = 20$. Kan $f(1) = 12$? Forklar.
\end{tasks}

\begin{sjekk}
\begin{tasks}(2)
  \task Finn vekstfaktoren ved $9\,\%$ vekst.
  \task Regn ut $f(3)$ når $f(x) = 7 \cdot 2^x$.
  \task! Er $f(x) = 90 \cdot 0{,}95^x$ voksende eller avtagende? Forklar.
\end{tasks}
\end{sjekk}

% ---------------------------------------------------------------------

\del{2}{Potensfunksjoner}

\begin{ordliste}
potensfunksjon & En funksjon $f(x) = a \cdot x^k$. Variabelen $x$ står i grunntallet. & \mittsprak\\
eksponent & Tallet $k$ i $x^k$. Det sier hvor mange ganger $x$ er faktor. & \mittsprak\\
proporsjonal & To størrelser der den ene dobles når den andre dobles: $y = a \cdot x$. & \mittsprak\\
modell & En funksjon som beskriver en virkelig situasjon omtrent, ikke nøyaktig. & \mittsprak\\
\end{ordliste}

\begin{mustkunne}
\textbf{Potensfunksjon:} $f(x) = a \cdot x^k$. Eksempler: $x^2$, $x^3$, $\sqrt{x} = x^{\tfrac12}$ og $\frac{1}{x} = x^{-1}$.
Når $k$ er en brøk eller negativ, bruker vi $x > 0$, som i 1.3.\par\smallskip
\textbf{Hva skjer når $x$ dobles?}
\[ f(2x) = a \cdot (2x)^k = a \cdot 2^k \cdot x^k = 2^k \cdot f(x) \]
Dobler vi $x$, ganges $f(x)$ med $2^k$. Ganger vi $x$ med $r$, ganges $f(x)$ med $r^k$.\par\smallskip
\textbf{Tre typer. Hva sier en tabell?}\par\smallskip
\renewcommand{\arraystretch}{1.4}
\begin{tabular}{@{}l l l@{}}
\textbf{Type} & \textbf{Funksjon} & \textbf{Mønster i tabellen}\\ \hline
Lineær & $ax + b$ & $x$ øker med 1: $y$ øker med et fast tall\\
Eksponentiell & $a \cdot b^x$ & $x$ øker med 1: $y$ ganges med en fast faktor\\
Potens & $a \cdot x^k$ & $x$ ganges med 2: $y$ ganges med en fast faktor\\
\end{tabular}
\end{mustkunne}

\begin{eksempel}[1 — dobling og tredobling]
Se på $f(x) = 3x^2$. Hva skjer med $f(x)$ når $x$ dobles? Når $x$ tredobles?

\begin{stegtabell}
$f(2x) = 3 \cdot (2x)^2$ & Sett inn $2x$ for $x$.\\
$= 3 \cdot 4 \cdot x^2 = 4 \cdot 3x^2 = 4 \cdot f(x)$ & Potens av produkt: $(2x)^2 = 4x^2$.\\
$f(3x) = 3 \cdot 9 \cdot x^2 = 9 \cdot f(x)$ & Samme regning: $3^2 = 9$.\\
\end{stegtabell}
Dobles $x$, blir $f(x)$ fire ganger så stor. Tredobles $x$, blir $f(x)$ ni ganger så stor.
\end{eksempel}

\begin{eksempel}[2 — hvilken type er tabellen?]
\begin{tabular}{@{}c|cccc@{}}
$x$ & 1 & 2 & 4 & 8\\ \hline
$y$ & 2 & 16 & 128 & 1024\\
\end{tabular}

\begin{stegtabell}
$\frac{16}{2} = 8$, $\frac{128}{16} = 8$, $\frac{1024}{128} = 8$ & $x$ dobles hver gang. $y$ ganges med 8 hver gang.\\
$2^k = 8$, så $k = 3$ & Potensfunksjon, fordi $x$ ganges med 2 og $y$ med en fast faktor.\\
$f(1) = a \cdot 1^3 = 2$, så $a = 2$ & Sett $x = 1$.\\
$f(x) = 2x^3$ & Kontroll: $f(4) = 2 \cdot 64 = 128$.\\
\end{stegtabell}
Tabellen er ikke lineær: $\frac{\Delta y}{\Delta x}$ er $14$, $56$ og $224$, og det er ikke konstant. Den er ikke eksponentiell: fra $x = 1$ til $2$ er forholdet 8, så vi ville fått $b = 8$. Fra $x = 2$ til $4$ er det to steg, og da skulle forholdet vært $8^2 = 64$. Men det er 8.
\end{eksempel}

\begin{vanligfeil}
\begin{itemize}[leftmargin=1.2em,itemsep=1pt,topsep=0pt]
  \item $2^x$ og $x^2$ er ikke samme funksjon. For $x = 3$ er $2^3 = 8$ og $3^2 = 9$. I $2^x$ er $x$ eksponent, i $x^2$ er $x$ grunntall.
  \item $x^{\tfrac12} = \frac{x}{2}$ gjelder ikke generelt. Det er kvadratroten, og for $x = 9$ er $9^{\tfrac12} = 3$, men $\frac{9}{2} = 4{,}5$.
  \item «Dobler vi $x$, dobles $f(x)$» gjelder ikke generelt. Det gjelder bare når $k = 1$. For $x^2$ firedobles $f(x)$, fordi $(2x)^2 = 4x^2$.
\end{itemize}
\end{vanligfeil}

\nivaahode{1}

\oppg{2.1}{1}{\uh}
Er funksjonen en eksponentialfunksjon eller en potensfunksjon?
\begin{tasks}(4)
  \task $f(x) = x^5$
  \task $g(x) = 5^x$
  \task $h(x) = \dfrac{4}{x}$
  \task $k(x) = 0{,}5^x$
\end{tasks}

\oppg{2.2}{1}{\uh}
Regn ut uten kalkulator.
\begin{tasks}(2)
  \task $f(x) = x^{\tfrac12}$. Finn $f(25)$.
  \task $g(x) = x^{-1}$. Finn $g(4)$.
  \task $h(x) = 2x^3$. Finn $h(3)$.
  \task $k(x) = x^{\tfrac32}$. Finn $k(9)$.
\end{tasks}

\oppg{2.3}{1}{\uh}
Se på $f(x) = 4x^3$. Hva skjer med $f(x)$ når $x$ dobles? Fullfør.
\[ f(2x) = 4 \cdot (2x)^3 = 4 \cdot \ \skrivlinje[12mm] \ \cdot x^3 = \ \skrivlinje[12mm] \ \cdot f(x) \]

\nivaahode{2}

\oppg{2.4}{2}{\uh}
Hva ganges $f(x)$ med når $x$ endres slik?
\begin{tasks}(3)
  \task $f(x) = 5x^3$, og $x$ tredobles.
  \task $g(x) = x^{\tfrac12}$, og $x$ firedobles.
  \task $h(x) = \dfrac{10}{x}$, og $x$ tredobles.
\end{tasks}

\oppg{2.5}{2}{\uh}
Hver tabell er lineær, eksponentiell eller en potensfunksjon.
\begin{center}
\begin{tabular}{@{}l|cccc@{}}
Tabell A: $x$ & 0 & 1 & 2 & 3\\
$y$ & 7 & 12 & 17 & 22\\
\end{tabular}\qquad
\begin{tabular}{@{}l|cccc@{}}
Tabell B: $x$ & 0 & 1 & 2 & 3\\
$y$ & 5 & 15 & 45 & 135\\
\end{tabular}\qquad
\begin{tabular}{@{}l|cccc@{}}
Tabell C: $x$ & 1 & 2 & 4 & 8\\
$y$ & 24 & 12 & 6 & 3\\
\end{tabular}
\end{center}
\begin{tasks}(1)
  \task Avgjør type for hver tabell. Begrunn med differanser eller forhold.
  \task Skriv funksjonsuttrykk for tabell B og tabell C.
\end{tasks}

\oppg{2.6}{2}{\uh}
Figuren viser grafene til $x^{\tfrac12}$, $x$ og $x^2$ for $x \ge 0$.
\begin{center}
\begin{tikzpicture}
\begin{axis}[width=6cm, height=5cm, axis lines=middle, xmin=0, xmax=4.4, ymin=0, ymax=4.6,
  xtick={1,2,3,4}, ytick={1,2,3,4}, tick label style={font=\scriptsize}, xlabel={$x$}, ylabel={$y$}, label style={font=\small}, grid=both, grid style={gray!25}]
\addplot[domain=0:2, thick, hoved, samples=50] {x^2};
\addplot[domain=0:4.2, thick, red!70!black, samples=2] {x};
\addplot[domain=0:4.2, thick, green!50!black, samples=60] {sqrt(x)};
\node[font=\small\bfseries, hoved, anchor=west] at (axis cs:1.95,3.7) {A};
\node[font=\small\bfseries, red!70!black, anchor=south] at (axis cs:3.85,3.9) {B};
\node[font=\small\bfseries, green!50!black, anchor=north] at (axis cs:4.1,1.9) {C};
\end{axis}
\end{tikzpicture}
\end{center}
\begin{tasks}(1)
  \task Hvilken graf hører til hvilken funksjon?
  \task For hvilke $x$ er $x^{\tfrac12}$ størst av de tre? For hvilke $x$ er $x^2$ størst?
  \task Hva er felles for de tre funksjonene når $x = 1$? Forklar hvorfor rekkefølgen bytter der.
\end{tasks}

\oppg{2.7}{2}{\uh}
Finn $k$ slik at $f(x) = x^k$ gir $f(4) = 8$. Finn \textbf{den første feilen}. Løs riktig.
\[ 4^k = 8 \ \longrightarrow\ (2^2)^k = 2^3 \ \longrightarrow\ 2^{k+2} = 2^3 \ \longrightarrow\ k + 2 = 3 \ \longrightarrow\ k = 1 \]

\nivaahode{3}

\oppg{2.8}{3}{\uh}
Er påstanden alltid, noen ganger eller aldri sann? Begrunn svaret.
Vi ser på $x > 0$ og $a > 0$.
\begin{tasks}(1)
  \task $f(x) = a \cdot x^k$ er voksende når $k > 0$.
  \task $f(x) = x^k$ er avtagende.
  \task Grafen til $f(x) = \frac{a}{x}$ skjærer $x$-aksen.
\end{tasks}

\oppg{2.9}{3}{\mh}
Bruk graftegner eller GeoGebra. Se på $2^x$ og $x^2$ for $x \ge 0$.
\begin{tasks}(1)
  \task For hvilke $x \ge 0$ er $2^x > x^2$?
  \task Se på $x$ fra 10 til 11. Hva skjer med hver verdi? Forklar hvorfor $2^x$ vinner for store $x$.
  \task Undersøk negative $x$. Hvor mange skjæringspunkter har grafene til sammen?
\end{tasks}

\oppg{2.10}{3}{\uh}
En funksjon $f$ går gjennom $(2, 20)$ og $(4, 80)$.
\begin{tasks}(1)
  \task Finn $a$ og $k$ i potensfunksjonen $f(x) = a \cdot x^k$.
  \task Finn $a$ og $b$ i eksponentialfunksjonen $g(x) = a \cdot b^x$ gjennom de samme punktene.
  \task Hvilket punkt ville avslørt typen? Begrunn uten å regne ut $f(6)$.
\end{tasks}

\begin{sjekk}
\begin{tasks}(2)
  \task Hva ganges $f(x) = x^3$ med når $x$ dobles?
  \task Regn ut $3^x$ og $x^3$ for $x = 2$.
  \task! Skriv $\dfrac{1}{x^2}$ som en potensfunksjon.
\end{tasks}
\end{sjekk}
